\input{../tex-inputs/latex-leseansicht-vorspann}

\section[Olga und Arthur Schnitzler an Gustav Schwarzkopf, 8. 7. 1909]{L04458 Olga und Arthur Schnitzler an Gustav Schwarzkopf, 8. 7. 1909}
\nopagebreak\mylabel{L04458v}
\rehead{ }\normalsize\beginnumbering\briefempfaengerindex{Schwarzkopf, Gustav@\textsc{Schwarzkopf, Gustav}!zzzSchnitzler, Arthur@\emph{von Arthur Schnitzler}!1909-07-081@{8. 7. 1909}|(be}\briefempfaengerindex{Schwarzkopf, Gustav@\textsc{Schwarzkopf, Gustav}!zzzSchnitzler, Olga@\emph{von Olga Schnitzler}!1909-07-081@{8. 7. 1909}|(be}
\toendnotes[C]{\smallbreak\pagebreak[2]}
\correspDesc{Versand  durch Olga Schnitzler, Arthur Schnitzler am 8. 7. 1909 in Edlach
\newline{}Übermittlung  am 9. 7. 1909 in Edlach
\newline{}Erhalt  durch Gustav Schwarzkopf im Zeitraum [9. 7. 1909
                  – 13. 7. 1909?] in Wien}\toendnotes[C]{\smallbreak}
\Standort{DLA, A:Schnitzler, HS.1985.1.1897.}
\physDesc{Bildpostkarte, 271 Zeichen
\newline{}Handschrift Olga Schnitzler: Bleistift, lateinische Kurrent
\newline{}Handschrift Arthur Schnitzler: Bleistift, deutsche Kurrent
\newline{}Versand: Stempel: »\nobreak{}\oindex{Edlach@\textbf{Edlach}|pwk}Edlach b. Reichenau in N.Oe, 9/7 09, 2–6 \textcolor{gray}{N}\nobreak{}«.  }\toendnotes[C]{\smallbreak}\pstart{}{\pb}Herrn\pend{}\pstart{}\textsc{Gustav Schwarzkopf}\pend{}\pstart{}Wien I\oindex{I., Innere Stadt@\textbf{I., Innere Stadt}|pw}\pend{}\pstart{}\textsc{Tiefer Graben 23.}\oindex{Wien@\textbf{Wien}!I., Innere Stadt@\textbf{I., Innere Stadt}!Tiefer Graben 23@\textbf{Tiefer Graben 23}|pw}\pend{}{\bigskip}
\pstart
           \noindent{}{\pb}\textcolor{gray}{\textbf{Edlach\oindex{Edlach@\textbf{Edlach}|pw} bei Reichenau\oindex{Reichenau an der Rax@\textbf{Reichenau an der Rax}|pw}}}\pend
           \vspace{1em}
\pstart
           \noindent{}{\pb}Lieber Herr Schwarzkopf, ich bin Ihnen sehr dankbar dafür, dass Sie
               den Heini\pwindex{Schnitzler, Heinrich 9.\,8.\,1902 Hinterbrühl – 12.\,7.\,1982 Wien@\textsc{Schnitzler, Heinrich} (9.\,8.\,1902 Hinterbrühl – 12.\,7.\,1982 Wien)|pw} besucht haben, ich \label{K_L04458-1v}\edtext{darf es ja leider nicht}{\lemma{\textnormal{\emph{darf es ja leider nicht}}}\Cendnote{\textnormal{Im \emph{Tagebuch}\pwindex{Schnitzler, Arthur 15. 5. 1862 Wien – 21. 10. 1931 ebd.@\textsc{Schnitzler, Arthur} (15. 5. 1862 Wien – 21. 10. 1931 ebd.)!Tagebuch@\strich\emph{Tagebuch}|pwk} vermerkte Arthur
                     Schnitzler in der Woche zuvor, dass der Sohn\pwindex{Schnitzler, Heinrich 9.\,8.\,1902 Hinterbrühl – 12.\,7.\,1982 Wien@\textsc{Schnitzler, Heinrich} (9.\,8.\,1902 Hinterbrühl – 12.\,7.\,1982 Wien)|pwkv} »wegen seines Hustens noch
                     in Wien, und nicht im Hotel wohnen« dürfe, siehe A. S.: \emph{Tagebuch}, 1. 7. 1909.}}}\label{K_L04458-1} tun.
               Hoffentlich sehen wir Sie bald heraussen! Indess herzliche Grüsse!\pend
           
\pstart
           \spacefill\mbox{Olga Schnitzler.}{\\[\baselineskip]}{[}hs. A. Schnitzler:{]} Ihr \spacefill\mbox{Arthur.}\pend
           \leftskip=0em{}
\pstart
           {[}hs. O. Schnitzler:{]} 8. Juli 09.\pend
           \selectlanguage{ngerman}\endnumbering\briefempfaengerindex{Schwarzkopf, Gustav@\textsc{Schwarzkopf, Gustav}!zzzSchnitzler, Arthur@\emph{von Arthur Schnitzler}!1909-07-081@{8. 7. 1909}|)be}\briefempfaengerindex{Schwarzkopf, Gustav@\textsc{Schwarzkopf, Gustav}!zzzSchnitzler, Olga@\emph{von Olga Schnitzler}!1909-07-081@{8. 7. 1909}|)be}\mylabel{L04458h}
\begin{anhang}
\end{anhang}\newcommand{\dateiname}{L04458}\newcommand{\titel}{Olga und Arthur Schnitzler an Gustav Schwarzkopf, 8. 7. 1909}\newcommand{\editorInnen}{Herausgegeben von Jahnke, SelmaMüller, Martin Anton}\input{../tex-inputs/latex-leseansicht-abspann}
